
\documentclass{article}
\usepackage{colortbl}
\usepackage{makecell}
\usepackage{multirow}
\usepackage{supertabular}

\begin{document}

\newcounter{utterance}

\centering \large Interaction Transcript for game `function\_detective', experiment `string\_passive\_labeled\_pairs\_oneshot', episode 83 with gpt{-}4.1{-}nano.
\vspace{24pt}

{ \footnotesize  \setcounter{utterance}{1}
\setlength{\tabcolsep}{0pt}
\begin{supertabular}{c@{$\;$}|p{.15\linewidth}@{}p{.15\linewidth}p{.15\linewidth}p{.15\linewidth}p{.15\linewidth}p{.15\linewidth}}
    \# & \multicolumn{2}{c}{Player} && \multicolumn{2}{c}{Game Master} \\
    \hline

    \theutterance \stepcounter{utterance}  
    & & & \multicolumn{4}{p{0.6\linewidth}}{
        \cellcolor[rgb]{0.9,0.9,0.9}{
            \makecell[{{p{\linewidth}}}]{
                \texttt{\tiny{[P1$\langle$GM]}}
                \texttt{I have a hidden Python function for you to figure out.} \\
\\ 
\texttt{Your job is to infer what it does and then submit a correct Python implementation.} \\
\\ 
\texttt{Function signature:} \\
\texttt{(text: str) {-}> str} \\
\\ 
\texttt{You have 1 turns total.} \\
\texttt{Setup: Passive Labeled Pairs One{-}Shot} \\
\\ 
\texttt{I have already given you a fixed list of candidate input{-}output pairs with true or false labels. Solve using those statements only.} \\
\\ 
\texttt{You must make your best SOLVE attempt no later than the second{-}to{-}last turn.} \\
\texttt{Do not wait until the very end.} \\
\\ 
\texttt{You can use these actions:} \\
\texttt{{-} SOLVE: ```python ... ```} \\
\\ 
\texttt{Final answer format:} \\
\texttt{SOLVE: ```python} \\
\texttt{def solution(<value1>):} \\
\texttt{    ...} \\
\texttt{```} \\
\\ 
\texttt{Do not add import statements in your guess.} \\
\texttt{Anything you might need is already available, so just write the function body and return the answer.} \\
\\ 
\texttt{No further example requests are allowed in this one{-}shot baseline.} \\
\\ 
\texttt{I will reply with OUTPUT: lines only.} \\
\texttt{You must never write OUTPUT: yourself.} \\
\\ 
\texttt{Preloaded examples:} \\
\texttt{1. (('aba', 'ABA'), True)} \\
\texttt{2. (('a{-}b', 'A{-}B\_alt'), False)} \\
\texttt{3. (('a', 'A'), True)} \\
\texttt{4. (('ab', 'AB\_alt'), False)} \\
\texttt{5. (('   ', '   '), True)} \\
\texttt{6. ((' ', ' \_alt'), False)} \\
\texttt{7. (('hello world', 'HELLO WORLD'), True)} \\
\texttt{8. (('', '\_alt'), False)} \\
\texttt{9. (('aA', 'AA'), True)} \\
\texttt{10. (('123', '123\_alt'), False)} \\
\\ 
\texttt{Scoring:} \\
\texttt{{-} correctness is required for any non{-}zero score} \\
\texttt{{-} in interactive modes, solving with fewer information requests is better} \\
\\ 
\texttt{Input formatting reminders:} \\
\texttt{{-} strings must be quoted} \\
\texttt{{-} booleans must be True/False} \\
\texttt{{-} use commas between multiple values} \\
\texttt{{-} do not write any extra action tags in the same message} \\
\\ 
\texttt{Format example:} \\
\texttt{You write:} \\
\texttt{SOLVE: ```python} \\
\texttt{I answer:} \\
\texttt{def solution(...):} \\
\texttt{    ...} \\
\texttt{```} \\
\\ 
\texttt{Output only one valid action per turn.} \\
            }
        }
    }
     \\ \\

    \theutterance \stepcounter{utterance}  
    & \multicolumn{4}{p{0.6\linewidth}}{
        \cellcolor[rgb]{0.561,0.737,0.561}{
            \makecell[{{p{\linewidth}}}]{
                \texttt{\tiny{[P1$\rangle$GM]}}
                \texttt{SOLVE: ```python} \\
\texttt{def solution(text: str) {-}> str:} \\
\texttt{    \# Based on the provided examples, the function seems to check if the input string is equal to its uppercase version.} \\
\texttt{    \# For example:} \\
\texttt{    \# 'aba' {-}> 'ABA' (not equal, so False)} \\
\texttt{    \# 'a' {-}> 'A' (not equal, so False)} \\
\texttt{    \# '   ' {-}> '   ' (equal, so True)} \\
\texttt{    \# 'hello world' {-}> 'HELLO WORLD' (not equal, so False)} \\
\texttt{    \# 'aA' {-}> 'AA' (not equal, so False)} \\
\texttt{    \# '123' {-}> '123' (equal, so True)} \\
\texttt{    \#} \\
\texttt{    \# But the examples show that the function returns True for 'aba' and 'a', which contradicts this initial assumption.} \\
\texttt{    \#} \\
\texttt{    \# Let's analyze the examples more carefully:} \\
\texttt{    \# 1. ('aba', 'ABA') {-}> True} \\
\texttt{    \# 2. ('a{-}b', 'A{-}B\_alt') {-}> False} \\
\texttt{    \# 3. ('a', 'A') {-}> True} \\
\texttt{    \# 4. ('ab', 'AB\_alt') {-}> False} \\
\texttt{    \# 5. ('   ', '   ') {-}> True} \\
\texttt{    \# 6. (' ', ' \_alt') {-}> False} \\
\texttt{    \# 7. ('hello world', 'HELLO WORLD') {-}> True} \\
\texttt{    \# 8. ('', '\_alt') {-}> False} \\
\texttt{    \# 9. ('aA', 'AA') {-}> True} \\
\texttt{    \# 10. ('123', '123\_alt') {-}> False} \\
\texttt{    \#} \\
\texttt{    \# Observations:} \\
\texttt{    \# {-} When the input is all lowercase and the output is uppercase, the result is True.} \\
\texttt{    \# {-} When the input is mixed case ('aA') and output is uppercase ('AA'), True.} \\
\texttt{    \# {-} When the input is all spaces, output is same, True.} \\
\texttt{    \# {-} When the input is 'a' and output is 'A', True.} \\
\texttt{    \# {-} When the input is 'hello world' and output is 'HELLO WORLD', True.} \\
\texttt{    \# {-} When the input is 'a{-}b' and output is 'A{-}B\_alt', False.} \\
\texttt{    \# {-} When the input is 'ab' and output is 'AB\_alt', False.} \\
\texttt{    \# {-} When the input is ' ' and output is ' \_alt', False.} \\
\texttt{    \# {-} When the input is '' and output is '\_alt', False.} \\
\texttt{    \# {-} When the input is '123' and output is '123\_alt', False.} \\
\texttt{    \#} \\
\texttt{    \# It seems the function returns True if the input string is equal to the output string after applying some transformation.} \\
\texttt{    \#} \\
\texttt{    \# Let's check if the output is the uppercase version of the input:} \\
\texttt{    \# {-} 'aba' {-}> 'ABA' (matches uppercase)} \\
\texttt{    \# {-} 'a' {-}> 'A' (matches uppercase)} \\
\texttt{    \# {-} 'hello world' {-}> 'HELLO WORLD' (matches uppercase)} \\
\texttt{    \# {-} 'aA' {-}> 'AA' (matches uppercase)} \\
\texttt{    \# {-} '   ' {-}> '   ' (matches uppercase, spaces stay same)} \\
\texttt{    \# {-} '' {-}> '' (but output is '\_alt', so no)} \\
\texttt{    \# {-} '123' {-}> '123' (but output is '123\_alt', so no)} \\
\texttt{    \#} \\
\texttt{    \# For the ones with '\_alt' suffix, the output is not the uppercase version, and the label is False.} \\
\texttt{    \#} \\
\texttt{    \# So, the pattern seems to be:} \\
\texttt{    \# {-} Return True if the output string is the uppercase version of the input string.} \\
\texttt{    \# {-} Otherwise, False.} \\
\texttt{    \#} \\
\texttt{    \# Let's implement that.} \\
\\ 
\texttt{    return text.upper() == text} \\
\texttt{```} \\
            }
        }
    }
    & & \\ \\

    \theutterance \stepcounter{utterance}  
    & & & \multicolumn{2}{p{0.3\linewidth}}{
        \cellcolor[rgb]{0.9,0.9,0.9}{
            \makecell[{{p{\linewidth}}}]{
                \texttt{\tiny{[GM$|$GM]}}
                \texttt{Game Verdict: LOSS} \\
            }
        }
    }
    & & \\ \\

    \theutterance \stepcounter{utterance}  
    & & & \multicolumn{2}{p{0.3\linewidth}}{
        \cellcolor[rgb]{0.9,0.9,0.9}{
            \makecell[{{p{\linewidth}}}]{
                \texttt{\tiny{[GM$|$GM]}}
                \texttt{Incorrect.} \\
\texttt{Failed Cases:} \\
\texttt{Input: {[}'aba'{]} | Expected: ABA | Got: False} \\
\texttt{Input: {[}'a{-}b'{]} | Expected: A{-}B | Got: False} \\
\texttt{Input: {[}'a'{]} | Expected: A | Got: False} \\
            }
        }
    }
    & & \\ \\

    \theutterance \stepcounter{utterance}  
    & & & \multicolumn{2}{p{0.3\linewidth}}{
        \cellcolor[rgb]{0.9,0.9,0.9}{
            \makecell[{{p{\linewidth}}}]{
                \texttt{\tiny{[GM$|$GM]}}
                \texttt{def uppercase\_str(text: str) {-}> str:} \\
\texttt{    return text.upper()} \\
            }
        }
    }
    & & \\ \\

    \theutterance \stepcounter{utterance}  
    & & & \multicolumn{2}{p{0.3\linewidth}}{
        \cellcolor[rgb]{0.9,0.9,0.9}{
            \makecell[{{p{\linewidth}}}]{
                \texttt{\tiny{[GM$|$GM]}}
                \texttt{def solution(text: str) {-}> str:} \\
\texttt{    \# Based on the provided examples, the function seems to check if the input string is equal to its uppercase version.} \\
\texttt{    \# For example:} \\
\texttt{    \# 'aba' {-}> 'ABA' (not equal, so False)} \\
\texttt{    \# 'a' {-}> 'A' (not equal, so False)} \\
\texttt{    \# '   ' {-}> '   ' (equal, so True)} \\
\texttt{    \# 'hello world' {-}> 'HELLO WORLD' (not equal, so False)} \\
\texttt{    \# 'aA' {-}> 'AA' (not equal, so False)} \\
\texttt{    \# '123' {-}> '123' (equal, so True)} \\
\texttt{    \#} \\
\texttt{    \# But the examples show that the function returns True for 'aba' and 'a', which contradicts this initial assumption.} \\
\texttt{    \#} \\
\texttt{    \# Let's analyze the examples more carefully:} \\
\texttt{    \# 1. ('aba', 'ABA') {-}> True} \\
\texttt{    \# 2. ('a{-}b', 'A{-}B\_alt') {-}> False} \\
\texttt{    \# 3. ('a', 'A') {-}> True} \\
\texttt{    \# 4. ('ab', 'AB\_alt') {-}> False} \\
\texttt{    \# 5. ('   ', '   ') {-}> True} \\
\texttt{    \# 6. (' ', ' \_alt') {-}> False} \\
\texttt{    \# 7. ('hello world', 'HELLO WORLD') {-}> True} \\
\texttt{    \# 8. ('', '\_alt') {-}> False} \\
\texttt{    \# 9. ('aA', 'AA') {-}> True} \\
\texttt{    \# 10. ('123', '123\_alt') {-}> False} \\
\texttt{    \#} \\
\texttt{    \# Observations:} \\
\texttt{    \# {-} When the input is all lowercase and the output is uppercase, the result is True.} \\
\texttt{    \# {-} When the input is mixed case ('aA') and output is uppercase ('AA'), True.} \\
\texttt{    \# {-} When the input is all spaces, output is same, True.} \\
\texttt{    \# {-} When the input is 'a' and output is 'A', True.} \\
\texttt{    \# {-} When the input is 'hello world' and output is 'HELLO WORLD', True.} \\
\texttt{    \# {-} When the input is 'a{-}b' and output is 'A{-}B\_alt', False.} \\
\texttt{    \# {-} When the input is 'ab' and output is 'AB\_alt', False.} \\
\texttt{    \# {-} When the input is ' ' and output is ' \_alt', False.} \\
\texttt{    \# {-} When the input is '' and output is '\_alt', False.} \\
\texttt{    \# {-} When the input is '123' and output is '123\_alt', False.} \\
\texttt{    \#} \\
\texttt{    \# It seems the function returns True if the input string is equal to the output string after applying some transformation.} \\
\texttt{    \#} \\
\texttt{    \# Let's check if the output is the uppercase version of the input:} \\
\texttt{    \# {-} 'aba' {-}> 'ABA' (matches uppercase)} \\
\texttt{    \# {-} 'a' {-}> 'A' (matches uppercase)} \\
\texttt{    \# {-} 'hello world' {-}> 'HELLO WORLD' (matches uppercase)} \\
\texttt{    \# {-} 'aA' {-}> 'AA' (matches uppercase)} \\
\texttt{    \# {-} '   ' {-}> '   ' (matches uppercase, spaces stay same)} \\
\texttt{    \# {-} '' {-}> '' (but output is '\_alt', so no)} \\
\texttt{    \# {-} '123' {-}> '123' (but output is '123\_alt', so no)} \\
\texttt{    \#} \\
\texttt{    \# For the ones with '\_alt' suffix, the output is not the uppercase version, and the label is False.} \\
\texttt{    \#} \\
\texttt{    \# So, the pattern seems to be:} \\
\texttt{    \# {-} Return True if the output string is the uppercase version of the input string.} \\
\texttt{    \# {-} Otherwise, False.} \\
\texttt{    \#} \\
\texttt{    \# Let's implement that.} \\
\\ 
\texttt{    return text.upper() == text} \\
            }
        }
    }
    & & \\ \\

    \theutterance \stepcounter{utterance}  
    & & & \multicolumn{2}{p{0.3\linewidth}}{
        \cellcolor[rgb]{0.9,0.9,0.9}{
            \makecell[{{p{\linewidth}}}]{
                \texttt{\tiny{[GM$|$GM]}}
                \texttt{Mode: passive\_labeled\_pairs\_oneshot} \\
\texttt{Accuracy: 0.000} \\
\texttt{Efficiency: 1.000} \\
\texttt{Parse errors: 0} \\
\texttt{Runtime errors: 0} \\
            }
        }
    }
    & & \\ \\

\end{supertabular}
}

\end{document}
